% Part: turing-machines
% Chapter: machines-computations
% Section: halting-states

\documentclass[../../../include/open-logic-section]{subfiles}

\begin{document}

\olfileid[fr]{tur}{mac}{hal}
\olsection{États d'arrêt}

\begin{explain}
Nous avons défini nos machines de telle sorte qu'elles ne
s'arrêtent que lorsqu'aucune instruction ne peut être exécutée.
Dans d'autres présentations usuelles des machines de Turing,
on distingue un \emph{état d'arrêt}~$h$, avec $h\in Q$.

L'idée est simple : lorsque la machine a terminé son travail
(elle a décidé d'accepter l'entrée ou fini d'écrire la sortie),
elle passe dans l'état~$h$, où elle s'arrête. Certaines machines
possèdent deux états d'arrêt, l'un pour accepter l'entrée et
l'autre pour la rejeter.
\end{explain}

\begin{ex}\emph{États d'arrêt}.
Pour éclairer cette notion, modifions d'abord la machine paire.
Au lieu de la laisser s'arrêter dans l'état~$q_0$ lorsque
l'entrée est paire, ajoutons une instruction qui la fait
passer dans un état d'arrêt.
\[
\begin{tikzpicture}[->,>=stealth',shorten >=1pt,auto,node distance=2.8cm,
                    semithick]
  \tikzstyle{every state}=[fill=none,draw=black,text=black]

  \node[initial, state]         (A)                     {$q_0$};
  \node[state]         (B) [right of=A] {$q_1$};
  \node[state]         (C) [below of=A] {$h$};

  \path (A) edge [bend left] node {\TMtrans{\TMstroke}{\TMstroke}{R}} (B)
  	    edge node {\TMtrans{\TMblank}{\TMblank}{N}} (C)
        (B) edge [loop above] node {\TMtrans{\TMblank}{\TMblank}{R}} (B)
            edge [bend left] node {\TMtrans{\TMstroke}{\TMstroke}{R}} (A);
\end{tikzpicture}
\]

Poussons l'exemple plus loin. Quand la machine constate que
l'entrée est impaire, elle ne s'arrête jamais. On peut la
modifier pour y ajouter un état de \emph{rejet}~$r$ : il suffit
de remplacer l'instruction qui la fait boucler par une
instruction qui la conduit dans cet état. Dans cette présentation à deux états d'arrêt, le rejet est distingué de l'acceptation par l'état final ; dans notre convention initiale, rejeter signifie ne pas s'arrêter.
\[
\begin{tikzpicture}[->,>=stealth',shorten >=1pt,auto,node distance=2.8cm,
                    semithick]
  \tikzstyle{every state}=[fill=none,draw=black,text=black]

  \node[initial,state]         (A)                     {$q_0$};
  \node[state]         (B) [right of=A] {$q_1$};
  \node[state]         (C) [below of=A] {$h$};
  \node[state]         (D) [below of=B] {$r$};

  \path (A) edge [bend left] node {\TMtrans{\TMstroke}{\TMstroke}{R}} (B)
  	    edge node {\TMtrans{\TMblank}{\TMblank}{N}} (C)
        (B) edge node {\TMtrans{\TMblank}{\TMblank}{N}} (D)
            edge [bend left] node {\TMtrans{\TMstroke}{\TMstroke}{R}} (A);
\end{tikzpicture}
\]
\end{ex}

\begin{explain}
Un état d'arrêt désigné peut être utile dans un cas comme
celui-ci : il rend explicite l'acceptation ou le rejet de
certaines entrées. Il ne donne cependant aucune puissance
de calcul supplémentaire. Notre convention d'arrêt, moins
chargée en états particuliers, a elle aussi des avantages.
La définition employée jusqu'ici dans ce chapitre rend
la démonstration de l'indécidabilité du \emph{problème de l'arrêt} intuitive
et facile à suivre. Nous la conserverons donc.
\end{explain}

\end{document}
